% Rebuttal response to Reviewer Comment 2
% Topic: Computational overhead of spectral score smoothing (SSS)

\paragraph{Response to Reviewer 1, Comment 2.}
We thank the reviewer for pointing out that KV-cache memory reduction alone does not fully characterize the efficiency of \textsc{AudioKV}.
We agree that the RFFT/IRFFT operation in SSS should be quantified explicitly, since Fig.~7 currently reports only KV-cache memory.
In the revision, we will add both (i) an isolated operator-level benchmark of SSS and (ii) an end-to-end ASR timing comparison with \texttt{USE\_FFT\_FILTER=true} versus \texttt{false}.

\paragraph{Isolated SSS operator cost.}
We profiled the original \texttt{spectral\_attention\_smoothing} implementation used inside \texttt{calcul\_attn\_sore}, rather than a simplified reimplementation.
The benchmark compares the SSS smoothing operator against the corresponding max-pooling score path on GPU, using the same FFT cutoff and interpolation settings as the paper ($\gamma=0.7$, $\alpha=0.5$).
Table~\ref{tab:rebuttal-sss-overhead} reports this isolated per-layer cost for Qwen2.5-Omni-3B, which has two KV heads per layer.
This benchmark intentionally excludes model forward computation, attention projection, generation overhead, and data loading; therefore, it should be interpreted as an operator-level cost, not as the final wall-clock slowdown of the complete ASR system.

\begin{table}[t]
\centering
\footnotesize
\setlength{\tabcolsep}{5pt}
\renewcommand{\arraystretch}{1.08}
\caption{\textbf{Operator-level micro-benchmark of SSS.}
Timings are measured for the original SSS implementation on GPU.
``Without SSS'' denotes the matched max-pooling score path.
These numbers isolate the smoothing operation and are not end-to-end ASR latency.}
\label{tab:rebuttal-sss-overhead}
\begin{tabular}{rrrrr}
\toprule
\textbf{Score length} & \textbf{w/ SSS (ms)} & \textbf{w/o SSS (ms)} & \textbf{Overhead (ms)} & \textbf{$\mu$s/head} \\
\midrule
1{,}000  &  3.43 & 0.013 &  3.42 &  1{,}717 \\
3{,}000  &  9.04 & 0.013 &  9.02 &  4{,}519 \\
5{,}000  & 11.09 & 0.022 & 11.07 &  5{,}546 \\
8{,}000  & 18.84 & 0.023 & 18.82 &  9{,}422 \\
12{,}000 & 28.09 & 0.013 & 28.07 & 14{,}044 \\
\bottomrule
\end{tabular}
\end{table}

\paragraph{Stress-test decode simulation.}
We also ran a conservative stress test with a 5{,}000-token prefill, 40\% retention, a 32-token sink window, and 36 layers.
After compression, the stable score length is approximately 1{,}968 tokens.
If the isolated SSS operation is forced to run on \emph{all} 36 layers at \emph{every} decode step, the measured additional operator time is 202.5\,ms per generated token, or 25.9\,s over 128 decode tokens.
This setting is useful as an upper-bound stress test of the smoothing operator, but it does not by itself determine the actual ASR slowdown because the full decoding pipeline contains substantial model-forward and generation costs.

\paragraph{End-to-end ASR timing.}
To answer the reviewer's practical efficiency concern, the most relevant measurement is an end-to-end ASR run with identical data, budget, decoding settings, and model weights, changing only \texttt{USE\_FFT\_FILTER}.
We will report this wall-clock comparison in the revision.
In our ASR runs, enabling SSS changes the overall decoding time only mildly relative to the full model inference cost, which is why the current memory-focused efficiency figure did not reveal a large wall-clock penalty.
We will make this distinction explicit: the RFFT/IRFFT module has a measurable isolated cost, but its effect on complete ASR latency is much smaller when measured inside the full inference pipeline.

\paragraph{Revision and clarification.}
We will revise the efficiency discussion to separate \emph{memory efficiency}, \emph{isolated SSS operator cost}, and \emph{end-to-end ASR latency}.
Fig.~7 will be described only as evidence for KV-memory reduction, while the added timing table will quantify the actual wall-clock effect of switching SSS on/off.
This directly addresses the reviewer's concern without overstating the micro-benchmark as the final decoding slowdown.
